\documentclass[../../main.tex]{subfiles}
\begin{document}

\paragraph{【新課程】}
{\bf [解]}
$\cos\theta = c$，$s = \sin\theta$とおく．円柱側面の点$P$は$P(c, s, z)$とおける．ただし$0 \leqq \theta < 2\pi$とする．すると，$D$は
\begin{align}
  0 \leqq z \leqq 1 - c + \sqrt{3}s
\end{align}
で与えられる．$z$の存在条件から
\begin{align}
  &0 \leqq 1 - c + \sqrt{3}s \nonumber\\
  &\Longleftrightarrow \cos\left(\theta+\dfrac{\pi}{3}\right) \leqq \dfrac{1}{2} \nonumber\\
  &\Longleftrightarrow 0 \leqq \theta \leqq \dfrac{4\pi}{3}
\end{align}
となる．

したがって，求める面積$S$は
\begin{align}
  S &= \int_0^{\frac{4\pi}{3}} (1 - c + \sqrt{3}s) \, d\theta \nonumber\\
  &= \Bigl[ \theta - s - \sqrt{3}c \Bigr]_0^{\frac{4\pi}{3}} \nonumber\\
  &= 2\sqrt{3} + \dfrac{4\pi}{3}
\end{align}
となる．$\cdots\text{（答）}$

\paragraph{【旧課程】}
{\bf [解]}
時刻$t$における点$P$の座標を$(x(t), 0)$とする．
条件より$x(0) = 0$であり，$P$は$x$軸上を正の向きに動くから，$t > 0$において$x(t) > 0$である．
$x \geqq 0$において$\dfrac{x^2+1}{x+1} > 0$であるから，時刻$t$において囲まれる図形の面積$S$は
\begin{align}
  S = \int_0^x \dfrac{u^2+1}{u+1} \, du
\end{align}
と表される．微分積分学の基本定理より
\begin{align}
  \dfrac{dS}{dx} = \dfrac{x^2+1}{x+1}
\end{align}
である．合成関数の微分法により
\begin{align}
  \dfrac{dS}{dt} = \dfrac{dS}{dx} \cdot \dfrac{dx}{dt} = \dfrac{x^2+1}{x+1} \dfrac{dx}{dt}
\end{align}
となる．題意より$P$が点$(x, 0)$を通過するときの変化率は$\dfrac{dS}{dt} = x^2+1$であるから，
\begin{align}
  \dfrac{x^2+1}{x+1} \dfrac{dx}{dt} = x^2+1
\end{align}
が成り立つ．$x^2+1 > 0$であるから両辺を$x^2+1$で割ると
\begin{align}
  \dfrac{1}{x+1} \dfrac{dx}{dt} = 1 \Longleftrightarrow \dfrac{dx}{dt} = x + 1
\end{align}
を得る．変数分離形として両辺を$t$で積分すると
\begin{align}
  \int \dfrac{1}{x+1} \, dx = \int 1 \, dt
\end{align}
より，$\log_e(x+1) = t + C$（$C$は積分定数）となる（$x \geqq 0$より$x+1 > 0$）．
$t = 0$のとき$x = 0$であるから，$\log_e 1 = 0 + C$より$C = 0$である．
したがって
\begin{align}
  \log_e(x+1) = t \Longleftrightarrow x+1 = e^t \Longleftrightarrow x = e^t - 1
\end{align}
である．
ここで，面積$S$の定積分を計算すると，
\begin{align}
  \dfrac{u^2+1}{u+1} = \dfrac{(u^2-1) + 2}{u+1} = u - 1 + \dfrac{2}{u+1}
\end{align}
であるから，
\begin{align}
  S &= \int_0^x \left( u - 1 + \dfrac{2}{u+1} \right) \, du \nonumber\\
  &= \left[ \dfrac{u^2}{2} - u + 2\log_e(u+1) \right]_0^x \nonumber\\
  &= \dfrac{x^2}{2} - x + 2\log_e(x+1)
\end{align}
となる．$x = e^t - 1$，および$\log_e(x+1) = t$を代入して整理すると，
\begin{align}
  S &= \dfrac{(e^t - 1)^2}{2} - (e^t - 1) + 2t \nonumber\\
  &= \dfrac{e^{2t} - 2e^t + 1}{2} - e^t + 1 + 2t \nonumber\\
  &= \dfrac{1}{2}e^{2t} - 2e^t + 2t + \dfrac{3}{2} \tag{答}
\end{align}
である．

\end{document}
